\setcounter{chapter}{4}
\chapter{Lebesgue測度}

\paragraph{本章の目標}
Lebesgue外測度 $\lambda^*$ とLebesgue内測度 $\lambda_*$ が一致する集合を
Lebesgue可測集合と定める．
そのうえで，Lebesgue可測集合全体が $\sigma$-加法族をなし，
Lebesgue測度が可算加法的な測度になることを確認する．
最後に，内外正則性とCarathéodory条件による判定法を述べる．

\section{Lebesgue可測集合とLebesgue測度}

以下では，
\[
Q_R:=[-R,R]^d\qquad(R>0)
\]
と書く．

\begin{definition}[Lebesgue可測集合とLebesgue測度]
集合 $A \subset \RR^d$ がLebesgue可測であるとは
\[
\lambda_*(A\cap Q_R)=\lambda^*(A\cap Q_R)
\qquad (R>0)
\]
が成り立つことをいう．
Lebesgue可測集合全体を
\[
\calL_d:=\{A\subset\RR^d\mid A \text{ はLebesgue可測}\}
\]
と書く．
$A\in\calL_d$ のとき，
\[
\lambda(A):=\lambda^*(A)
\]
を $A$ のLebesgue測度という．
\end{definition}

% \begin{definition}[Lebesgue零集合]
% 集合 $N\subset\RR^d$ が
% \[
% \lambda^*(N)=0
% \]
% を満たすとき，$N$ をLebesgue零集合という．
% \end{definition}

\begin{theorem}[Lebesgue零集合]
Lebesgue零集合はLebesgue可測であり，Lebesgue測度は $0$ である．
\end{theorem}
\begin{proof}
\cref{prop:lebesgue-inner-basic}（Lebesgue内測度は外測度以下であること）より
\[
0\le \lambda_*(N)\le \lambda^*(N)
\]
である．
$\lambda^*(N)=0$ なら
\[
\lambda_*(N)=\lambda^*(N)=0
\]
だから，$N$ はLebesgue可測で $\lambda(N)=0$ である．
\end{proof}

\begin{corollary}[可算集合]
一点集合，可算集合，$\QQ^d$，$\ZZ^d$ はLebesgue可測であり，測度 $0$ をもつ．
%特に1次元では $\QQ$ と $\ZZ$ はLebesgue測度 $0$ である．
\end{corollary}
\begin{proof}
一点集合および可算集合のLebesgue外測度は
\cref{thm:countable-outer-measure-zero}（可算集合は外測度 $0$ であること）より $0$ である．
したがって前定理より従う．
\end{proof}

\begin{theorem}[Jordan可測集合] \label{thm:lebesgue-jordan-measurable}
Jordan可測集合 $A\subset\RR^d$ はLebesgue可測であり，
\[
\lambda(A)=m_J(A)
\]
が成り立つ．
\end{theorem}
\begin{proof}
Jordan可測性より
\[
m_{J,*}(A)=m_J^*(A)=m_J(A)
\]
である．
\cref{thm:jordan-inner-lebesgue-inner}（Jordan内測度はLebesgue内測度以下であること），
\cref{prop:lebesgue-inner-basic}（Lebesgue内測度は外測度以下であること），
\cref{prop:louter-le-jouter}（Lebesgue外測度はJordan外測度以下であること）を順に用いると
\[
m_J(A)
=m_{J,*}(A)
\le \lambda_*(A)
\le \lambda^*(A)
\le m_J^*(A)
=m_J(A)
\]
となる．
したがってすべて等号であり，特に
\[
\lambda_*(A)=\lambda^*(A)=m_J(A)
\]
である．
よって $A$ はLebesgue可測であり，
\[
\lambda(A)=m_J(A)
\]
である．
\end{proof}

\begin{example}
\[
\QQ\cap[0,1]
\]
は可算集合なのでLebesgue可測であり，
\[
\lambda_1(\QQ\cap[0,1])=0
\]
である．
一方，$[0,1]$ はJordan可測で $m_J([0,1])=1$ だから
\[
\lambda_1([0,1])=1
\]
である．
よって，後で示す有限加法性より
\[
\lambda_1([0,1]\setminus\QQ)=1
\]
となる．
\end{example}

\section{同値な条件}

\begin{theorem}[判定法] \label{thm:lebesgue-measurable-equivalent}
$A\subset\RR^d$ とする．
次は同値である．
\begin{enumerate}
    \item $A$ はLebesgue可測である
    \item $A$ はCarathéodory条件を満たす
    \item 任意の $R>0$ と $\eps>0$ に対し，開集合 $G$ とコンパクト集合 $K$ が存在して
    \[
    K\subset A\cap[-R,R]^d\subset G,
    \qquad
    \lambda^*(G\setminus K)<\eps
    \]
    を満たす
\end{enumerate}
さらに $\lambda^*(A)<\infty$ のときは，これらは次の条件とも同値である．
\begin{enumerate}
    \item[(4)] 任意の $\eps>0$ に対して開集合 $G$ とコンパクト集合 $K$ が存在して
    \[
    K\subset A\subset G,
    \qquad
    \lambda^*(G\setminus K)<\eps
    \]
    を満たす
\end{enumerate}
\end{theorem}
\begin{proof}
(1) と (2) の同値は，Lebesgue流の可測性
$\lambda_*(A\cap Q_R)=\lambda^*(A\cap Q_R)$ と
Carathéodory条件を直接照合して得る．
有限外測度の集合については，第4章の開集合を用いた内測度表示と
第3章の開集合による外正則性により，内外一致とCarathéodory条件が同値であることを確認できる．
一般の集合については，$A\cap Q_R$ に局所化してから
$R\to\infty$ とすればよい．

(1) から (3) を示す．
\[
A_R:=A\cap[-R,R]^d
\]
とおく．
$A_R$ はLebesgue可測で，しかも有限測度である．
内正則性より，コンパクト集合 $K\subset A_R$ で
\[
\lambda(A_R)-\lambda(K)<\frac{\eps}{2}
\]
を満たすものが取れる．
外正則性より，開集合 $G\supset A_R$ で
\[
\lambda(G)-\lambda(A_R)<\frac{\eps}{2}
\]
を満たすものが取れる．
すると $K\subset G$ であり，
\[
\lambda^*(G\setminus K)
=
\lambda(G\setminus K)
=
\lambda(G)-\lambda(K)
<\eps
\]
である．

(3) から (1) を示す．
各 $R>0$ について $A_R=A\cap[-R,R]^d$ とおく．
(3) により，任意の $\eps>0$ に対して
\[
K\subset A_R\subset G,
\qquad
\lambda^*(G\setminus K)<\eps
\]
を満たすコンパクト集合 $K$ と開集合 $G$ が取れる．
このとき $K$ と $G$ はLebesgue可測であり，
\[
\lambda^*(A_R)-\lambda_*(A_R)
\le
\lambda(G)-\lambda(K)
=
\lambda(G\setminus K)
<\eps
\]
である．
$\eps>0$ は任意だから $A_R$ はLebesgue可測である．
したがって
\[
A=\bigcup_{n=1}^{\infty}\bigl(A\cap[-n,n]^d\bigr)
\]
はLebesgue可測である．

最後に $\lambda^*(A)<\infty$ と仮定する．
(1) から (4) は，内正則性と外正則性を直接 $A$ に適用すればよい．
(4) から (1) は，上の (3) から (1) の証明と同じ議論を $A$ そのものに適用すれば従う．
\end{proof}

\begin{remark}
条件 (4) の形は有限測度の集合に対する判定法である．
例えば $A=\RR^d$ の場合，どのコンパクト集合 $K$ を取っても
\[
\RR^d\setminus K
\]
は無限の測度をもつので，(4) は成り立たない．
一般のLebesgue可測集合に対しては，条件 (3) のように有界な窓で切った局所版を用いる．
\end{remark}


\section{Lebesgue可測集合の性質}

\begin{theorem}[Lebesgue可測集合の $\sigma$-加法族性]
\label{thm:lebesgue-measurable-sigma-algebra}
Lebesgue可測集合全体 $\calL_d$ は $\RR^d$ 上の $\sigma$-加法族である．
すなわち，次が成り立つ．
\begin{enumerate}
    \item $\emptyset,\RR^d\in\calL_d$
    \item $A\in\calL_d$ ならば $\RR^d\setminus A\in\calL_d$
    \item $\{A_n\}_{n=1}^{\infty}\subset\calL_d$ ならば
    \[
    \bigcup_{n=1}^{\infty}A_n\in\calL_d
    \]
\end{enumerate}
\end{theorem}
\begin{proof}
\cref{thm:lebesgue-measurable-equivalent}（Lebesgue可測性とCarathéodory条件の同値性）より，
Lebesgue可測集合について示す代わりに，
Lebesgue外測度 $\lambda^*$ に関するCarathéodory条件を満たす集合全体が
$\sigma$-加法族であることを示せばよい．

以下，Carathéodory条件を満たす集合全体を $\calM$ と書く．
\cref{thm:lebesgue-outer-measure}（Lebesgue外測度は外測度であること）より，
$\lambda^*$ は単調性と可算劣加法性を満たす．

まず $\emptyset,\RR^d\in\calM$ である．
実際，任意の $E\subset\RR^d$ に対して
\[
\lambda^*(E)
=
\lambda^*(E\cap\RR^d)+\lambda^*(E\setminus\RR^d)
=
\lambda^*(E)+\lambda^*(\emptyset)
\]
であり，同様に
\[
\lambda^*(E)
=
\lambda^*(E\cap\emptyset)+\lambda^*(E\setminus\emptyset)
=
\lambda^*(\emptyset)+\lambda^*(E)
\]
である．

次に，$\calM$ が補集合で閉じていることを示す．
$A\in\calM$ とする．
任意の $E\subset\RR^d$ に対して
\[
E\cap A^c=E\setminus A,
\qquad
E\setminus A^c=E\cap A
\]
だから，$A$ のCarathéodory条件より
\[
\lambda^*(E)
=
\lambda^*(E\cap A)+\lambda^*(E\setminus A)
=
\lambda^*(E\cap A^c)+\lambda^*(E\setminus A^c).
\]
したがって $A^c\in\calM$ である．

次に，$\calM$ が有限合併で閉じていることを示す．
$A,B\in\calM$ とする．
任意の $E\subset\RR^d$ を取る．
まず $A$ のCarathéodory条件を $E$ に適用し，さらに
$B$ のCarathéodory条件を $E\setminus A$ に適用すると
\[
\lambda^*(E)
=
\lambda^*(E\cap A)+\lambda^*(E\setminus A)
=
\lambda^*(E\cap A)+\lambda^*(E\cap(B\setminus A))
  +\lambda^*(E\setminus(A\cup B)).
\]
一方，$A$ のCarathéodory条件を $E\cap(A\cup B)$ に適用すると
\[
\lambda^*(E\cap(A\cup B))
=
\lambda^*(E\cap A)+\lambda^*(E\cap(B\setminus A)).
\]
よって
\[
\lambda^*(E)
=
\lambda^*(E\cap(A\cup B))+\lambda^*(E\setminus(A\cup B)).
\]
したがって $A\cup B\in\calM$ である．
帰納法により，$\calM$ は有限合併で閉じている．
補集合と有限合併で閉じているので，有限交叉と差集合でも閉じている．

最後に，$\calM$ が可算合併で閉じていることを示す．
$A_n\in\calM$ $(n\ge1)$ とし，
\[
C:=\bigcup_{n=1}^{\infty}A_n
\]
とおく．
差集合で閉じていることから
\[
B_1:=A_1,\qquad
B_n:=A_n\setminus\bigcup_{k=1}^{n-1}A_k\quad(n\ge2)
\]
はすべて $\calM$ に属し，互いに素であり，
\[
C=\bigsqcup_{n=1}^{\infty}B_n
\]
である．
また
\[
C_N:=\bigsqcup_{n=1}^{N}B_n
\]
とおくと，有限合併で閉じていることから $C_N\in\calM$ である．
互いに素な有限個の集合に対してCarathéodory条件を繰り返し適用すると，
任意の $E\subset\RR^d$ について
\[
\lambda^*(E\cap C_N)
=
\sum_{n=1}^{N}\lambda^*(E\cap B_n)
\]
である．
さらに $C_N\in\calM$ だから
\[
\lambda^*(E)
=
\lambda^*(E\cap C_N)+\lambda^*(E\setminus C_N)
=
\sum_{n=1}^{N}\lambda^*(E\cap B_n)+\lambda^*(E\setminus C_N).
\]
ここで $E\setminus C_N\supset E\setminus C$ なので，単調性より
\[
\lambda^*(E)
\ge
\sum_{n=1}^{N}\lambda^*(E\cap B_n)+\lambda^*(E\setminus C).
\]
$N\to\infty$ とすると
\[
\lambda^*(E)
\ge
\sum_{n=1}^{\infty}\lambda^*(E\cap B_n)+\lambda^*(E\setminus C).
\]
一方，可算劣加法性より
\[
\lambda^*(E\cap C)
=
\lambda^*\left(\bigcup_{n=1}^{\infty}(E\cap B_n)\right)
\le
\sum_{n=1}^{\infty}\lambda^*(E\cap B_n)
\]
であるから
\[
\lambda^*(E)
\ge
\lambda^*(E\cap C)+\lambda^*(E\setminus C)
\]
を得る．
逆向きの不等式
\[
\lambda^*(E)
\le
\lambda^*(E\cap C)+\lambda^*(E\setminus C)
\]
は，$E=(E\cap C)\cup(E\setminus C)$ と可算劣加法性から従う．
したがって
\[
\lambda^*(E)
=
\lambda^*(E\cap C)+\lambda^*(E\setminus C)
\]
であり，$C\in\calM$ である．

以上より，Carathéodory条件を満たす集合全体 $\calM$ は
$\sigma$-加法族である．
\cref{thm:lebesgue-measurable-equivalent} により $\calM=\calL_d$ だから，
$\calL_d$ も $\sigma$-加法族である．
\end{proof}

\begin{corollary}[集合演算]
\label{cor:lebesgue-measurable-set-operations}
$A,B\in\calL_d$ ならば
\[
A\cup B,\qquad A\cap B,\qquad A\setminus B
\]
はいずれもLebesgue可測である．
また $\{A_n\}_{n=1}^{\infty}\subset\calL_d$ ならば
\[
\bigcap_{n=1}^{\infty}A_n
\]
もLebesgue可測である．
\end{corollary}
\begin{proof}
有限合併は可算合併の特別な場合である．
また
\[
A\cap B=(A^c\cup B^c)^c,
\qquad
A\setminus B=A\cap B^c
\]
であり，可算共通部分は
\[
\bigcap_{n=1}^{\infty}A_n
=
\left(\bigcup_{n=1}^{\infty}A_n^c\right)^c
\]
と書ける．
\cref{thm:lebesgue-measurable-sigma-algebra} を用いれば，これらはいずれもLebesgue可測である．
\end{proof}

\begin{theorem}[開集合と閉集合]
\label{thm:open-closed-lebesgue-measurable}
開集合と閉集合はLebesgue可測である．
\end{theorem}
\begin{proof}
端点がすべて有理数である有界開区間
\[
\prod_{i=1}^{d}(a_i,b_i)
\qquad(a_i,b_i\in\QQ,\ a_i<b_i)
\]
の全体は可算であり，$\RR^d$ の開集合の基底をなす．
したがって任意の開集合 $G\subset\RR^d$ は，そのような有界開区間の可算和として書ける．

有界開区間はJordan可測であるから，
\cref{thm:lebesgue-jordan-measurable}（Jordan可測集合はLebesgue可測であること）よりLebesgue可測である．
\cref{thm:lebesgue-measurable-sigma-algebra} により可算和もLebesgue可測なので，$G$ はLebesgue可測である．
閉集合は開集合の補集合だから，同じく \cref{thm:lebesgue-measurable-sigma-algebra} よりLebesgue可測である．
\end{proof}

\section{Lebesgue測度の性質}

\begin{lemma}[有限加法性]
\label{lem:lebesgue-measure-finite-additive}
$A,B\in\calL_d$ が互いに素ならば
\[
\lambda(A\cup B)=\lambda(A)+\lambda(B)
\]
が成り立つ．
\end{lemma}
\begin{proof}
\cref{thm:lebesgue-measurable-equivalent} より，$A$ はCarathéodory条件を満たす．
$E=A\cup B$ としてその条件を適用すると，$A\cap B=\emptyset$ だから
\[
\lambda^*(A\cup B)
=
\lambda^*((A\cup B)\cap A)+\lambda^*((A\cup B)\setminus A)
=
\lambda^*(A)+\lambda^*(B)
\]
である．
\cref{cor:lebesgue-measurable-set-operations} より $A\cup B$ もLebesgue可測なので，
これは
\[
\lambda(A\cup B)=\lambda(A)+\lambda(B)
\]
を意味する．
\end{proof}

\begin{theorem}[可算加法性]
\label{thm:lebesgue-measure-countably-additive}
互いに素なLebesgue可測集合列 $\{A_n\}_{n=1}^{\infty}$ に対して
\[
\lambda\left(\bigsqcup_{n=1}^{\infty}A_n\right)
=
\sum_{n=1}^{\infty}\lambda(A_n)
\]
が成り立つ．
\end{theorem}
\begin{proof}
まず \cref{lem:lebesgue-measure-finite-additive} を繰り返すと，任意の $N$ について
\[
\lambda\left(\bigsqcup_{n=1}^{N}A_n\right)
=
\sum_{n=1}^{N}\lambda(A_n)
\]
である．
また
\[
\bigsqcup_{n=1}^{N}A_n
\subset
\bigsqcup_{n=1}^{\infty}A_n
\]
だから，\cref{thm:lebesgue-outer-measure}（Lebesgue外測度の単調性）より
\[
\sum_{n=1}^{N}\lambda(A_n)
\le
\lambda\left(\bigsqcup_{n=1}^{\infty}A_n\right)
\]
である．
$N\to\infty$ とすれば
\[
\sum_{n=1}^{\infty}\lambda(A_n)
\le
\lambda\left(\bigsqcup_{n=1}^{\infty}A_n\right)
\]
を得る．
逆向きの不等式は，\cref{thm:lebesgue-outer-measure}（Lebesgue外測度の可算劣加法性）より
\[
\lambda\left(\bigsqcup_{n=1}^{\infty}A_n\right)
=
\lambda^*\left(\bigsqcup_{n=1}^{\infty}A_n\right)
\le
\sum_{n=1}^{\infty}\lambda^*(A_n)
=
\sum_{n=1}^{\infty}\lambda(A_n)
\]
である．
\end{proof}

\begin{theorem}[基本性質]
\label{thm:lebesgue-measure-basic}
Lebesgue測度について次が成り立つ．
\begin{enumerate}
    \item 非負性：
    \[
    0\le \lambda(A)\le\infty
    \qquad(A\in\calL_d)
    \]
    かつ $\lambda(\emptyset)=0$ である
    \item 単調性：$A,B\in\calL_d$ かつ $A\subset B$ ならば
    \[
    \lambda(A)\le \lambda(B)
    \]
    \item 劣加法性：$\{A_n\}_{n=1}^{\infty}\subset\calL_d$ ならば
    \[
    \lambda\left(\bigcup_{n=1}^{\infty}A_n\right)
    \le
    \sum_{n=1}^{\infty}\lambda(A_n)
    \]
\end{enumerate}
\end{theorem}
\begin{proof}
\begin{enumerate}
    \item Lebesgue測度はLebesgue外測度の制限である．
    \cref{thm:lebesgue-outer-measure} より，Lebesgue外測度は $[0,\infty]$ に値をもち，
    \[
    \lambda(\emptyset)=\lambda^*(\emptyset)=0
    \]
    である．

    \item $A\subset B$ ならば，\cref{thm:lebesgue-outer-measure} の単調性より
    \[
    \lambda(A)=\lambda^*(A)\le \lambda^*(B)=\lambda(B)
    \]
    である．

    \item \cref{thm:lebesgue-measurable-sigma-algebra} より $\bigcup_n A_n$ はLebesgue可測である．
    したがって \cref{thm:lebesgue-outer-measure} の可算劣加法性より
    \[
    \lambda\left(\bigcup_{n=1}^{\infty}A_n\right)
    =
    \lambda^*\left(\bigcup_{n=1}^{\infty}A_n\right)
    \le
    \sum_{n=1}^{\infty}\lambda^*(A_n)
    =
    \sum_{n=1}^{\infty}\lambda(A_n)
    \]
    を得る．
\end{enumerate}
\end{proof}

\begin{theorem}[下からの連続性]
\label{thm:lebesgue-measure-continuity-from-below}
Lebesgue可測集合の単調増大列
\[
A_1\subset A_2\subset\cdots
\]
に対して
\[
\lambda\left(\bigcup_{n=1}^{\infty}A_n\right)
=
\lim_{n\to\infty}\lambda(A_n)
\]
が成り立つ．
\end{theorem}
\begin{proof}
\[
B_1:=A_1,\qquad
B_n:=A_n\setminus A_{n-1}\quad(n\ge2)
\]
とおく．
\cref{cor:lebesgue-measurable-set-operations} より各 $B_n$ はLebesgue可測であり，定義から互いに素である．
また
\[
A_N=\bigsqcup_{n=1}^{N}B_n,
\qquad
\bigcup_{n=1}^{\infty}A_n=\bigsqcup_{n=1}^{\infty}B_n
\]
である．
\cref{lem:lebesgue-measure-finite-additive} と \cref{thm:lebesgue-measure-countably-additive} より
\[
\lambda(A_N)=\sum_{n=1}^{N}\lambda(B_n),
\qquad
\lambda\left(\bigcup_{n=1}^{\infty}A_n\right)
=
\sum_{n=1}^{\infty}\lambda(B_n).
\]
$N\to\infty$ とすれば主張を得る．
\end{proof}

\begin{theorem}[上からの連続性]
\label{thm:lebesgue-measure-continuity-from-above}
Lebesgue可測集合の単調減少列
\[
A_1\supset A_2\supset\cdots
\]
に対して，$\lambda(A_1)<\infty$ ならば
\[
\lambda\left(\bigcap_{n=1}^{\infty}A_n\right)
=
\lim_{n\to\infty}\lambda(A_n)
\]
が成り立つ．
\end{theorem}
\begin{proof}
\[
B_n:=A_1\setminus A_n
\]
とおく．
すると
\[
B_1\subset B_2\subset\cdots,
\qquad
\bigcup_{n=1}^{\infty}B_n
=
A_1\setminus\bigcap_{n=1}^{\infty}A_n
\]
である．
\cref{thm:lebesgue-measure-continuity-from-below} より
\[
\lambda\left(A_1\setminus\bigcap_{n=1}^{\infty}A_n\right)
=
\lim_{n\to\infty}\lambda(A_1\setminus A_n).
\]
一方，$A_n\subset A_1$ かつ $\lambda(A_1)<\infty$ なので，
\cref{lem:lebesgue-measure-finite-additive} より
\[
\lambda(A_1\setminus A_n)=\lambda(A_1)-\lambda(A_n)
\]
である．
同様に
\[
\lambda\left(A_1\setminus\bigcap_{n=1}^{\infty}A_n\right)
=
\lambda(A_1)-\lambda\left(\bigcap_{n=1}^{\infty}A_n\right).
\]
これらを合わせれば主張が従う．
\end{proof}

\begin{proposition}[平行移動不変性]
\label{prop:lebesgue-measure-translation-invariant}
$A\in\calL_d$ と $c\in\RR^d$ に対して
\[
A+c:=\{x+c\mid x\in A\}
\]
もLebesgue可測であり，
\[
\lambda(A+c)=\lambda(A)
\]
が成り立つ．
\end{proposition}
\begin{proof}
\cref{thm:lebesgue-measurable-equivalent} より，$A$ はCarathéodory条件を満たす．
任意の $E\subset\RR^d$ に対して，\cref{prop:lebesgue-outer-translation-invariant}
（Lebesgue外測度は平行移動で不変であること）を用いると
\begin{align*}
\lambda^*(E)
&=
\lambda^*(E-c)\\
&=
\lambda^*((E-c)\cap A)+\lambda^*((E-c)\setminus A)\\
&=
\lambda^*(E\cap(A+c))+\lambda^*(E\setminus(A+c)).
\end{align*}
したがって $A+c$ もCarathéodory条件を満たす．
再び \cref{thm:lebesgue-measurable-equivalent} より，$A+c$ はLebesgue可測である．
さらに \cref{prop:lebesgue-outer-translation-invariant} より
\[
\lambda(A+c)=\lambda^*(A+c)=\lambda^*(A)=\lambda(A)
\]
である．
\end{proof}
